% =========================================================================
%  Ejercicio 2b - Ensayo INDIVIDUAL sobre el articulo
%  Responsable: Miguel Ángel Jiménez Ramírez
%
%  NOTA: Este ensayo es INDIVIDUAL, uno por cada integrante del equipo.
% =========================================================================

\subsection*{Ensayo individual — Miguel Ángel Jiménez Ramírez}

\section*{Introducción}
La cuarta revolución industrial está construida sobre una economía digital donde la información es el recurso más importante y dinámico. El problema principal que aborda la lectura trata sobre la mercantilización masiva de nuestros datos personales, ya que generamos información sin parar al usar dispositivos electrónicos y redes sociales, convirtiéndonos sin darnos cuenta en los activos de las grandes empresas tecnológicas. Desde el punto de vista tecnológico, esto es muy importante porque muestra una gran desigualdad de poder entre las personas y las instituciones que guardan y procesan todos estos datos. Creo que este tema es muy importante porque las bases de datos corporativas no solo almacenan registros sencillos, sino que modelan y aprovechan la vida digital de millones de personas sin que exista una conciencia clara sobre la privacidad, el control de la información y los derechos humanos que están en juego al gestionar bases de datos a gran escala.

\section*{Desarrollo}
De acuerdo con el artículo, la gran cantidad de información que producimos se rige por las leyes del mercado, funcionando bajo la oferta y la demanda dentro de un entorno dominado por el Big Data y el Internet de las Cosas (IoT). La idea principal del autor señala que la recolección automática y rápida de nuestros datos debilita la confianza en el mercado digital, transformando a las personas en simples ceros y unos guardados en servidores corporativos que usan perfiles de comportamiento para ganar más dinero con la publicidad dirigida. Esta situación se conecta sobre todo en lo referente a cómo se estructuran los esquemas de almacenamiento, cómo se organizan grandes volúmenes de registros y qué modelos de datos relacionales y no relacionales se necesitan para procesar transacciones complejas. Así como un Sistema Gestor de Bases de Datos (SGBD) organiza los datos para mantenerlos seguros, ordenados e independientes, la recolección masiva que menciona la lectura muestra cómo las empresas construyen grandes espacios de almacenamiento centralizados o distribuidos donde los datos personales como la ubicación, el historial de navegación y los mensajes sirven como base para tomar decisiones mediante algoritmos.

A partir de esto, considero que la supuesta gratuidad de los servicios de Internet es en realidad una trampa comercial que esconde un intercambio desigual, donde pagamos con nuestra propia privacidad e identidad digital. No está bien que las empresas ganen muchísimo dinero aprovechando los contenidos que suben los usuarios y las opiniones que dejamos en las plataformas, mientras que nosotros, los dueños de esa información, no tenemos herramientas reales para controlar, revisar o monetizar lo que dejamos en la red. Este problema se ve claramente en casos reales que menciona la lectura, como los experimentos donde investigadores en España subastaron datos personales o los estudios sobre cuánto vale la información digital, los cuales demuestran que las personas valoran más sus datos financieros y su ubicación que sus búsquedas de todos los días. Desde un punto de vista práctico, esto nos enseña que hace falta diseñar bases de datos enfocadas en las personas, como los Almacenes de Datos Personales (PDS) o plataformas descentralizadas como The Hub of All Things (HAT), las cuales permiten que cada usuario tenga el control y la responsabilidad sobre su propia información mediante un consentimiento claro y consciente.

Me parece que el objetivo de la lectura es hacernos entender el valor real que tiene nuestra información y las consecuencias éticas de que las empresas recopilen tantos datos, buscando que aprendamos a defender nuestra privacidad. El tema central del artículo gira en torno a la economía de los datos personales en la época del Big Data y el IoT, el cual se acompaña de otros temas secundarios como el derecho a la privacidad, la desigualdad de poder con las empresas, las leyes internacionales como el Reglamento General de Protección de Datos (GDPR) de Europa y los nuevos modelos para administrar la información de forma descentralizada. Estos temas se relacionan con el tema principal porque establecen las reglas y límites para manejar los datos. Para mí, esta relación es muy importante; un buen profesional no solo debe saber hacer consultas en SQL, normalizar tablas o administrar motores de bases de datos, sino que también debe cumplir con las leyes de privacidad, aplicar controles de acceso seguros y diseñar estructuras que protejan los datos sensibles contra robos o malos usos.

\section*{Conclusión}
En conclusión, se mantiene la idea inicial de que el uso comercial de nuestros datos personales en la red exige un cambio hacia la transparencia y el empowerment de los usuarios, para que el control de la información no quede únicamente en manos de las empresas. Lo más valioso de la lectura es que nos ayuda a ver que los datos no sonneutrales y que cada interacción digital es un recurso valioso que requiere nuevas formas de gestión centradas en las personas, apoyadas por tecnologías como los almacenes personales y mejores leyes. Todo esto abre nuevos retos para los profesionales del área, como crear sistemas que permitan cumplir con el derecho a la portabilidad y al olvido, diseñar bases de datos que protejan la privacidad desde el inicio, y encontrar un equilibrio entre la innovación tecnológica y el respeto a los derechos humanos en la era digital.

% Extension: entre 1 y 2 cuartillas por integrante.
%
% Estructura obligatoria (en parrafos, sin listas):
% - INTRODUCCION: presenta brevemente el tema, por que es importante, y
%   plantea tu postura o idea principal.
% - DESARROLLO: resume los argumentos del autor, relaciona con otros
%   conocimientos/lecturas/ejemplos, expone tus propios argumentos (a favor
%   o en contra) con razones, indica el objetivo que planteo el autor y como
%   se relaciona con la materia, y la tematica central y temas laterales y su
%   relacion con tu practica profesional.
% - CONCLUSION: reafirma tu postura, que aprendiste o la aportacion mas
%   valiosa, y si el tema abre preguntas o retos futuros.
% - FORMATO: parrafos completos, conectores logicos, ortografia y gramatica
%   cuidadas.

% Referencia del articulo (APA) - completar los datos:
% \cite{datavalue}